\section{Introduction}
\label{sec:introduction}

An agent runtime performs deterministic control work between model and tool
calls. It parses a typed outcome, advances a state machine, checks policy and
budget state, selects a route, and emits the next effect. Each transition is
small beside inference, but a service with many concurrent sessions executes
them continuously.

This control path is already a datacenter concern. A recent production and
open-source characterization found that agentic workflows repeatedly cross the
CPU--GPU boundary, place host orchestration on the critical path, and create
bursty CPU demand \cite{yang2026agenticarchitecture}. That result establishes
the pressure. It leaves a more specific systems question open: when can the
deterministic transition itself be grouped and executed profitably on a GPU?

Similarity across agents is not enough. A GPU route needs a cohort above its
hardware and runtime crossover, all members must share executable semantics,
and they must become ready before their launch deadlines. A decision returned
to the host after every step also incurs copy, synchronization, branch, and
redispatch overhead. We call the interface between cohort supply and
observation placement the \emph{ready-cohort boundary}.

For each route and hardware/runtime configuration, let $K$ begin a measured
safe suffix: every tested cohort from $K$ through a declared maximum clears a
chosen admissible baseline. A workload supplies events with release times,
launch deadlines, and declared grouping keys. The central workload question is
how much same-group work can be packed above $K$. The corresponding mechanism
question is how much of the control chain can proceed without exposing an
intermediate decision to the host.

We answer these questions in two experiments. A public agent-trace replay
compares fixed-window eligibility with an exact offline packing optimum and a
local upper bound. A separate CUDA study holds the state transition and route
bodies fixed while changing where one GPU-computed decision is observed. An
earlier device-launch design provides a negative control. The experiments are
kept numerically separate because the trace threshold $K=256$ is swept rather
than measured for the resident-policy source, and the 32-epoch mechanism
horizon is not inferred from the trace.

The paper makes four contributions:

\begin{enumerate}[label=\arabic*.]
  \item \textbf{A ready-cohort formulation.} We separate the hardware
  threshold $K$, fixed-partition share $F$, exact offline share $\pstar$, local
  bound $U$, and online achieved share $A$. Each quantity has one source and a
  defined inference boundary.
  \item \textbf{An exact opportunity instrument.} Under zero service,
  unlimited capacity, and equal relative launch deadlines, a specialized
  dynamic program computes $\pstar$. This gives a reproducible workload
  measurement without claiming algorithmic priority.
  \item \textbf{Trace evidence for hidden cohort supply.} In the frozen
  primary replay, exact packing raises eligible share from
  \TraceFixedPct\% to \TraceExactPct\% and recovers
  \TraceGapClosurePct\% of the fixed-window alignment gap. The full grid also
  identifies regimes with no qualifying cohort.
  \item \textbf{Cross-provider mechanism evidence.} Across four named GPU
  placements and \ResidentCells{} cells, retaining the tested decision on
  device is \ResidentRatioMin$\times$ to \ResidentRatioMax$\times$ faster
  than the matched host-mediated GPU path. A fixed nested graph loses in all
  \NativeCells{} cells across five placements, ruling out device launch alone
  as the explanation.
\end{enumerate}

These experiments resolve two prerequisite questions for a GPU control plane:
whether enough work can coexist under a declared grouping, and whether one
device-side decision avoids measurable host-mediated overhead. The next system
must join them with finite service and CPU fallback, then measure $A$, CPU use,
raw tail latency, exact effects, task utility, and inference interference.

\begin{figure*}[t]
  \centering
  \begin{tikzpicture}[
    node distance=0.75cm and 0.50cm,
    font=\sffamily\small,
    evidence/.style={draw=RCBlue!85!black,rounded corners=3pt,fill=RCBlue!6,
      text width=3.18cm,minimum height=1.12cm,align=center},
    next/.style={draw=RCCharcoal!70,rounded corners=3pt,fill=RCPaper,
      text width=3.18cm,minimum height=1.12cm,align=center,dashed},
    arrow/.style={-{Latex[length=2.2mm]},thick,draw=RCCharcoal!65}
  ]
    \node[evidence] (k) {\textbf{Threshold input $K$}\\swept in the trace study};
    \node[evidence,right=of k] (opp) {\textbf{Cohort supply}\\fixed $F$, exact $\pstar$, upper $U$};
    \node[next,right=of opp] (a) {\textbf{Online realization $A$}\\finite queues, service, fallback};
    \node[next,right=of a] (benefit) {\textbf{Service outcome}\\CPU, raw P99, utility, cost};
    \draw[arrow] (k) -- (opp);
    \draw[arrow] (opp) -- (a);
    \draw[arrow] (a) -- (benefit);
    \node[evidence,below=0.55cm of opp] (res)
      {\textbf{Observation placement}\\host round trip versus device decision};
    \draw[arrow] (res) -| (a);
  \end{tikzpicture}
  \caption{Evidence map. Solid boxes are the controlled or measured components
  in this paper; dashed boxes define the joined online experiment and its
  service-level outcomes. The shared notation connects the studies without
  treating the swept trace threshold as the resident source's measured
  crossover.}
  \label{fig:claim-boundary}
\end{figure*}
